% Gerado por regressao/06_tabelas_latex.R (projeto Modelagem). Não editar à mão.
\begin{sidewaystable}
\caption{Robustez R10 — sem limpeza (outliers de despesa e receita original)}\label{tab:robustez_R10}
\centering
\footnotesize
\setlength{\tabcolsep}{3pt}
\begin{tabular}{@{}>{\raggedright\arraybackslash}p{\dimexpr\linewidth-594pt\relax}>{\centering\arraybackslash}p{60pt}>{\centering\arraybackslash}p{60pt}>{\centering\arraybackslash}p{60pt}>{\centering\arraybackslash}p{60pt}>{\centering\arraybackslash}p{60pt}>{\centering\arraybackslash}p{60pt}>{\centering\arraybackslash}p{60pt}>{\centering\arraybackslash}p{60pt}>{\centering\arraybackslash}p{60pt}@{}}
\hline
\textbf{Variável} & \textbf{\hspace{0pt}Despesa total} & \textbf{\hspace{0pt}Saúde e Saneamento} & \textbf{\hspace{0pt}Educação e Cultura} & \textbf{\hspace{0pt}Habitação e Urbanismo} & \textbf{\hspace{0pt}Assistência Social} & \textbf{\hspace{0pt}Transporte} & \textbf{\hspace{0pt}Administração} & \textbf{\hspace{0pt}Agricultura} & \textbf{\hspace{0pt}Comunicações} \\ \hline
\multicolumn{10}{@{}l}{\textit{Modelo A (ciclo eleitoral)}} \\
Ano eleitoral & 284,240\textsuperscript{***} & 50,861\textsuperscript{***} & 50,146\textsuperscript{*} & 129,104\textsuperscript{***} & 14,228\textsuperscript{***} & 41,860\textsuperscript{***} & 9,846 & 1,220 & $-$0,749\textsuperscript{***} \\
 & (104,513) & (11,317) & (29,730) & (32,316) & (4,336) & (10,609) & (11,475) & (2,021) & (0,192) \\
Ano pré-eleitoral & 165,024\textsuperscript{*} & 2,216 & 78,867\textsuperscript{**} & 34,709\textsuperscript{**} & 9,534\textsuperscript{**} & 8,504 & 10,550 & 5,111 & 0,662\textsuperscript{***} \\
 & (95,972) & (12,822) & (38,886) & (16,941) & (4,694) & (7,758) & (8,369) & (3,668) & (0,145) \\
\hline
Observações & 70.855 & 70.698 & 70.723 & 69.741 & 70.658 & 54.551 & 70.730 & 64.625 & 12.382 \\
Municípios & 5.568 & 5.568 & 5.568 & 5.564 & 5.568 & 5.205 & 5.568 & 5.436 & 1.897 \\
R\textsuperscript{2} & 0,837 & 0,759 & 0,712 & 0,302 & 0,384 & 0,191 & 0,275 & 0,174 & 0,031 \\[6pt]
\multicolumn{10}{@{}l}{\textit{Modelo B (coalizões)}} \\
Prefeito na coalizão & 32,328\textsuperscript{***} & 3,784 & $-$0,761 & 4,797\textsuperscript{*} & 0,701 & 2,574 & 7,846\textsuperscript{***} & $-$0,514 & $-$0,091 \\
do governador & (6,457) & (2,387) & (3,107) & (2,803) & (0,723) & (2,654) & (2,776) & (1,049) & (0,304) \\
Prefeito na coalizão & $-$11,411 & 9,958\textsuperscript{***} & $-$1,930 & $-$0,591 & 1,044 & $-$12,312\textsuperscript{***} & 0,389 & $-$5,064\textsuperscript{***} & 0,392 \\
do presidente & (7,944) & (3,380) & (3,886) & (4,222) & (0,997) & (3,144) & (3,935) & (1,410) & (0,405) \\
\hline
Observações & 70.855 & 70.698 & 70.723 & 69.741 & 70.658 & 54.551 & 70.730 & 64.625 & 12.382 \\
Municípios & 5.568 & 5.568 & 5.568 & 5.564 & 5.568 & 5.205 & 5.568 & 5.436 & 1.897 \\
R\textsuperscript{2} & 0,578 & 0,431 & 0,344 & 0,110 & 0,163 & 0,093 & 0,138 & 0,071 & 0,023 \\
\hline
\end{tabular}
\fonte{Elaboração própria.}
\nota{Modelo A: efeito fixo de município, erros-padrão de Driscoll-Kraay. Modelo B: efeitos fixos de município e de ano, erros-padrão agrupados por município. Erros-padrão entre parênteses. \textsuperscript{***}p$<$0,01; \textsuperscript{**}p$<$0,05; \textsuperscript{*}p$<$0,1.}
\end{sidewaystable}
